
The solver is essentially a function that performs the mapping
\begin{equation*}
    \{\theta, L,n,\mathcal{M},\mathbf{x}_0\} \mapsto \mathbf{x}_{\text{sol}},
\end{equation*}
where $\theta$, $L$, and $n$ are known and $\mathbf{x}_{\text{sol}}$ is the proposed solution from the solver. The set of measurements, $\mathcal{M}$, is constructed by choosing some inlet measurements, $h_0^{(j)}$ and $q_0^{(j)}$, and simulating the outlet measurements, $h_L^{(j)}$ and $q_L^{(j)}$, by using the true leak parameters, $\mathbf{x^*}$, and the forward pass defined by \cref{eq:forward_pass}.
\begin{equation*}
    \begin{bmatrix}
        h_L^{(j)}\\
        q_L^{(j)}
    \end{bmatrix}
    = F_{fp}(\mathbf{x}, h_0^{(j)},q_0^{(j)}).
\end{equation*}
For the results in this section, $n$ sets of measurements were used, which yield a square system. Various strategies can be used for choosing the initial guess, $\mathbf{x}_0$.

\section{Discovery of Multiple Solutions}
\label{sec:result_multiple_solutions}

From running the solver on arbitrary leak configurations of two and three leaks, it was discovered that for most configurations, the solver would converge to both the true solution, $\mathbf{x}^*$, and to some other leak configuration, $\mathbf{x}^\dagger$. From extensive testing of arbitrary configurations, most leak configurations yield two solutions that solve the residual, $r(\mathbf{x})=0$. To further investigate the two solutions, the input-output relationship for both configurations was compared.

Let $\varepsilon$ be an error metric defined by
\begin{equation*}
    \varepsilon(h_0,q_0) = \|F_{fp}(\mathbf{x}^*\,h_0,q_0)-F_{fp}(\mathbf{x}^\dagger,h_0,q_0)\|.
\end{equation*}
To compare the two leak configurations, $\varepsilon$ is computed for a range of $h_0$ and $q_0$ values.



\section{Simulation Results}
% Present how I measure the performance of the solver. How many iterations, how are true leak parameters picked, how is initial x picked, how many runs, which solutions are counted, how is error rate calculated... etc.

An extensive test framework is designed to test the performance of the proposed solution method. The goal is to provide empirical results on the solver's performance across a wide range of leak configurations and pipe characteristics. For a given number of leaks, $n$, the test framework consists of running many test cases, where each test case consists of:
\begin{enumerate}
    \item Generate random leak parameters ($z_i$, $c_i$) and pipe resistance coefficient ($\theta$).
    \item Simulate $n$ sufficient sets of measurements ($\mu^{(1)},\dots,\mu^{(n)}$), creating a square system ($n=m$).
    \item Solve for the leak parameters using different strategies for choosing the initial value, $x_0$.
    \item Save results.
\end{enumerate}

By randomly generating test cases and running many different cases, the entire test suite will cover a wide range of configurations and provide a good picture of the methods' performance on simulated measurements. For simplicity, the pipe is normalized to unit length ($L=1$) and kept constant across test cases.\\

The leak locations, $z_1 < z_2 < \dots < z_n$, are randomly sampled from the uniform distribution,
\begin{equation}
\label{eq:zi_sampled}
    z_i \sim \mathcal{U}(\delta,1-\delta) \quad \text{s.t.} \quad z_{i+1}-z_i > \delta,
\end{equation}
where $\delta$ defines a small spacing that ensures a suitable configuration. For all test cases in this study, the spacing is constant and fixed at $\delta=0.05$.\\

The leak coefficients and pipe resistance coefficient are randomly sampled according to
\begin{equation}
\label{eq:ci_theta_sample}
    c_i \sim \mathcal{U}(0.01, \ 0.20), \qquad \theta \sim \mathcal{U}(0.1, \ 0.5).
\end{equation}\\

The set of measurements, $\mathcal{M}$, is generated from a constructed set of inlet heads, $h_0^{(j)}$, and randomly sampling $q_0^{(j)}$ until a sufficient pressure drop is obtained. The inlet heads are equally spaced in a defined interval and generated by
\begin{equation}
\label{eq:h0_gen}
    h_0^{(j)} = h_0^{\min} + (j-1)\dfrac{h_0^{\max}-h_0^{\min}}{n-1}.
\end{equation}
For the test cases in this study, the interval is fixed as $h_0^{\min}=5,\ h_0^{\max}=20$. The inlet flow is randomly sampled according to
\begin{equation}
\label{eq:q0_dist}
    q_0^{(j)} \sim \mathcal{U}(0.1, \ 20),
\end{equation}
and the outlet head and flow are simulated by the forward pass defined by \cref{eq:forward_pass}.
The inlet flow is resampled until a sufficient pressure drop, defined by
\begin{equation}
\label{eq:pressure_drop}
    \Delta h = 1-\frac{h_L^{(j)}}{h_0^{(j)}},
\end{equation}
occurs. For the test cases in this study, a pressure drop of 30\%-80\% ($\Delta h\in [0.3,0.8]$) is considered sufficient. The generation of the set of measurements is summarized in \autoref{alg:generate_measurements}.

\begin{algorithm}[ht]
\caption{Generation of Measurements}
\label{alg:generate_measurements}

\DontPrintSemicolon
\SetKwFunction{func}{GenerateMeasurements}
\SetKwFunction{linspace}{Linspace}
\SetKwFunction{return}{Return}
\SetKwProg{Fn}{Function}{:}{}
\SetKwInput{KwNote}{\normalfont Note}
\Fn{\func{$x,n$}}{
    $h_0 \gets$ \linspace{$h_0^{\min}, h_0^{\max},n$} \tcp*[r]{\cref{eq:h0_gen}}
    $\mathcal{M} \gets$ [ ] \;
    \For{$j \gets 1$ \KwTo n}{
        \Repeat{$0.3 < \Delta h < 0.8$}{
            $q_0^{(j)} \gets \mathcal{U}(0.1,20)$ \tcp*[r]{\cref{eq:q0_dist}}
            $h_L^{(j)},\ q_L^{(j)} \gets F_{fp}(x, h_0^{(j)},q_0^{(j)})$ \tcp*[r]{\cref{eq:forward_pass}}
            $\Delta h \gets 1-h_L^{(j)}/h_0^{(j)}$ \tcp*[r]{\cref{eq:pressure_drop}}
        }
        $\mathcal{M}$.append($h_0^{(j)},\ q_0^{(j)},\ h_L^{(j)},\ q_L^{(j)}$) \;
    }
    \return $\mathcal{M}$ \; 
}
\end{algorithm}

\subsection{Observability metric}

To model the detectability of an isolated leak, introduce a visibility threshold $c_{\mathrm{vis}}>0$ and define
\[
V_i=\frac{c_i}{c_i+c_{\mathrm{vis}}}.
\]
Hence, $V_i\approx 1$ for clearly visible leaks and $V_i\approx 0$ for very small leaks.

To model the loss of distinguishability caused by neighboring leaks, introduce a spatial resolution length $\ell>0$. The influence of a neighboring leak should increase when it is both close and large relative to $c_i$. This motivates the masking term
\[
M_i=
\exp\!\left(-\frac{L_i}{\ell}\right)\frac{c_{i-1}}{c_i+c_{i-1}}
+
\exp\!\left(-\frac{L_{i+1}}{\ell}\right)\frac{c_{i+1}}{c_i+c_{i+1}}.
\]
The exponential factors penalize small separation distances, while the amplitude ratios capture the fact that a small leak near a larger one is harder to identify.

The observability of leak $i$ is then defined as
\[
O_i = V_i\,\exp(-\beta M_i),
\]
where $\beta>0$ controls how strongly neighboring leaks reduce observability. Finally, a global observability metric is obtained by combining average and worst-case performance,
\[
O_{\mathrm{mean}}=\frac{1}{n}\sum_{i=1}^n O_i,\qquad
O_{\min}=\min_{1\le i\le n} O_i,\qquad
O_{\mathrm{global}}=\sqrt{O_{\mathrm{mean}}\,O_{\min}}.
\]
This construction yields a dimensionless score in $(0,1]$, which is maximized when leaks are sufficiently separated and have comparable magnitudes, and decreases when leaks are clustered or strongly imbalanced in size.





\subsection{Two Leaks}

\begin{table}[h!]
\centering
\begin{tabular}{c|c|cc|cc}
\hline
\multicolumn{2}{c|}{} & \multicolumn{4}{c}{\textbf{Parameters}} \\
\hline
\multicolumn{2}{c|}{\textbf{Case}} & $L_1$ & $L_2$ & $c_1$ & $c_2$ \\
\hline

\multirow{2}{*}{A}
& $P^{*}$        & 5.97   & 148.85 & $1.60\times10^{-4}$ & $2.30\times10^{-4}$ \\
& $\delta P^{*}$ & 0.03\% & 0.00\% & 0.00\%              & 0.00\%              \\
\hline

\multirow{2}{*}{B}
& $P^{*}$        & 41.05  & 78.89  & $2.20\times10^{-4}$ & $1.15\times10^{-4}$ \\
& $\delta P^{*}$ & 0.02\% & 0.01\% & 0.02\%              & 0.04\%              \\
\hline

\multirow{2}{*}{C}
& $P^{*}$        & 63.01  & 34.88  & $2.09\times10^{-4}$ & $2.09\times10^{-4}$ \\
& $\delta P^{*}$ & 0.06\% & 0.01\% & 0.21\%              & 0.21\%              \\
\hline

\multirow{2}{*}{D}
& $P^{*}$        & 5.97   & 57.04  & $1.50\times10^{-4}$ & $1.00\times10^{-4}$ \\
& $\delta P^{*}$ & 2.10\% & 0.07\% & 0.42\%              & 0.63\%              \\
\hline

\multirow{2}{*}{E}
& $P^{*}$        & 97.89  & 56.93  & $2.10\times10^{-4}$ & $1.10\times10^{-4}$ \\
& $\delta P^{*}$ & 0.00\% & 0.00\% & 0.00\%              & 0.01\%              \\
\hline
\end{tabular}
\caption{Error values for $L$ and $c$ across cases.}
\end{table}

\subsection{Three Leaks}

\section{Issues Solving Four Leaks}

Show that converging solutions are scattered all over the place.

Further analysed in \autoref{chap:results_issues_gt_3}